\renewcommand{\indexname}{INDICE ANALITICO}
\begin{theindex}

  \item \textbf{A}
  \indexspace
  \item Aberrazione (Angolo di) 91
  \item Alnico V 255
  \item Ammettenza 171
  \item \textsc{Amp\`ere} (amp\`ere, A: Unit\`a) 2, 74, 137
    \subitem (Teorema di) 125
  \item Angolo solido 8
  \item Antenna 202
    \subitem (Potenza irraggiata da una) 205
  \item Approssimazione di dipolo 203
  \item Armoniche (Funzioni) 35
  \item Autoinduzione 152

  \indexspace
  \item \textbf{B}
  \indexspace
  \item \textsc{Barkhausen} (Effetto) 253
  \item \textsc{Biot} -- \textsc{Savart} (Legge di) 127
  \item \textsc{Bohr} (Magnetone di) 244
    \subitem (Raggio di) 41
  \item \textsc{Bradley} 91

  \indexspace
  \item \textbf{C}
  \indexspace
  \item Caduta di potenziale 69
  \item Cammino libero medio 66
  \item Campo
    \subitem di radiazione 204
    \subitem elettrico (definizione) 4, 108
    \subitem -- di una distribuzione sferica 19
    \subitem -- cilindrica 19--20
    \subitem -- lineare 20
    \subitem -- piana 20$\div$22
    \subitem $\vec H$ 240
    \subitem irrotazionale 23
    \subitem magnetico 74--75, 111
    \subitem -- di saturazione 252, 253
    \subitem -- di una lamina di corrente 112, 134
    \subitem -- di una spira 128
    \subitem -- di un filo 134
    \subitem -- di un solenoide 129$\div$131, 134$\div$135
    \subitem -- di un toro 135
    \subitem -- rotante 182
    \subitem residuo 254
    \subitem solenoidale 124
  \item Capacit\`a
    \subitem (Coefficienti di) 47
    \subitem di un condensatore 51, 215, 225--226
    \subitem di un conduttore 50
  \item Cariche immagine 53
  \item Cariche magnetiche (sorgenti di $\vec H$) 249--250
  \item Cavit\`a in un dielettrico 227
  \item Cavo coassiale 155, 157--158
  \item CGS elettrostatico (Sistema) 2--3
  \item Ciclo di isteresi 254
    \subitem (Lavoro nel) 259
  \item Circuitazione 22
    \subitem per unit\`a di area 123
  \item Circuito
    \subitem antirisonante 172
    \subitem equivalente 69
    \subitem primario (o inducente) 150, 160
    \subitem risonante 172
    \subitem $LC$ 168
    \subitem $RC$ 71, 168
    \subitem $RL$ 158
    \subitem $RLC$ 166
    \subitem secondario (o indotto) 150, 160
  \item \textsc{Clausius} (v.\ \textsc{Mossotti})
  \item Coefficienti
    \subitem di autoinduzione 152
    \subitem di mutua induzione 150
    \subitem di capacit\`a 47
    \subitem di potenziale 47
  \item Coercitiva (Forza) 254
  \item Composizione delle velocit\`a 80, 88, 90--91, 94$\div$96
  \item Condensatore 51
    \subitem (Capacit\`a di un) 51, 215, 225--226
    \subitem (Carica e scarica del) 71$\div$73
  \item Conducibilit\`a 63
  \item Coniugazione di carica 18, 140
  \item Cono futuro, luce, passato 91
  \item Conservazione della carica 60--61, 125
  \item Conservazione del flusso nelle spire superconduttrici 188
  \item Contrazione delle lunghezze 89--90, 94
  \item Coordinate sferiche (o polari) 8
  \item Corrente
    \subitem alternata 61
    \subitem -- (Circuiti in) 169$\div$172
    \subitem continua (o stazionaria) 61
    \subitem (Densit\`a di) 59
    \subitem di magnetizzazione 239--240, 248
    \subitem -- in un magnete 252
    \subitem di polarizzazione 239
    \subitem di spostamento 146, 162
    \subitem indotta 148
    \subitem (Intensit\`a di) 59
    \subitem (Lamina di) 115, 134
    \subitem lentamente variabile 61, 146
    \subitem superficiale 112
  \item Costante dielettrica 215, 223
  \item Costante di \textsc{Faraday} 71
    \subitem di \textsc{Planck} 41, 212, 244
  \item \textsc{Coulomb} (Legge di) 1
    \subitem (coulomb, C: Unit\`a) 2, 3
  \item Covarianza della leggi fisiche 79
  \item \textsc{Curie} (Temperatura di) 241
  \item Curva di magnetizzazione 253

  \indexspace
  \item \textbf{D}
  \indexspace
  \item \textsc{d'Alembert} (Equazione di) 197
  \item Decadimento esponenziale 81--82
  \item Densit\`a
    \subitem di carica (v.\ Distribuzioni)
    \subitem di corrente 59
    \subitem di corrente di energia 189
    \subitem di energia 33, 189
    \subitem di magnetizzazione 240, 242, 248
    \subitem di polarizzazione 217, 238--239
    \subitem di quantit\`a di moto del campo 191
  \item Diamagnetismo 241, 245
  \item Difetto di massa 105
  \item Diffrazione 209$\div$211
  \item Dilatazione dei tempi 82, 89, 93--94
  \item Dipolo (Approssimazione di) 203
  \item Dipolo elettrico (v.\ anche Momento di dipolo) 18, 26--27
    \subitem (Forza su un) 232
    \subitem (Momento su un) 219
  \item Dipolo magnetico (v.\ Momento magnetico)
  \item \textsc{Dirichlet} (Problema di) 43
    \subitem Discontinuit\`a
    \subitem del campo elettrico 29
    \subitem del campo magnetico 131
  \item Distribuzioni di carica
    \subitem di volume 7
    \subitem di superficie 9
    \subitem lineari 9
  \item Divergenza 11
    \subitem (Teorema della) 12
  \item Domini di \textsc{Weiss} 252
  \item Dominio di armonicit\`a 35
  \item Doppio strato 41
  \item \textsc{Drude} -- \textsc{Lorentz} (Modello di) 65

  \indexspace
  \item \textbf{E}
  \indexspace
  \item Effetti di bordo 57--58, 236
  \item Effetto
    \subitem \textsc{Barkhausen} 253
    \subitem fotoelettrico 41
    \subitem \textsc{Joule} 64
    \subitem \textsc{Hall} 140--141
    \subitem pelle (o \textsc{Kelvin}) 63, 154--155
    \subitem termoionico 41
  \item \textsc{Einstein} 77
    \subitem (Relazione di) \ $E=h\nu$ \ 212
  \item Elettrone 2, 3
    \subitem (Raggio classico) 41, 194
    \subitem (Teoria di \textsc{Abraham} -- \textsc{Lorentz}) 194
  \item Elettronvolt (eV: Unit\`a) 24
  \item Energia (v.\ anche Densit\`a)
    \subitem del campo elettromagnetico 184--185
    \subitem del campo magnetico 156
    \subitem di un condensatore 53, 229--2230
    \subitem elettrostatica 31
    \subitem -- di un dielettrico 229
    \subitem libera 231, 258--259
    \subitem magnetica 258
    \subitem relativistica 102$\div$105
  \item Equazione
    \subitem delle onde 196
    \subitem di conservazione della carica 61, 125
    \subitem di continuit\`a per la corrente 61
    \subitem -- per l'energia 189
    \subitem di \textsc{d'Alembert} 197
    \subitem di \textsc{Laplace} 35
    \subitem di \textsc{Poisson} 28
  \item Equazioni costitutive 240
  \item Equazioni di Maxwell 121
    \subitem nel vuoto 196
    \subitem nella materia 240
  \item Equivalenza massa--energia 103$\div$105
  \item Extracorrente di apertura 160, 169

  \indexspace
  \item \textbf{F}
  \indexspace
  \item \textsc{Farad} (farad, F: Unit\`a) 50
  \item \textsc{Faraday} (Costante di) 71
    \subitem (Gabbia di) 45
    \subitem -- \textsc{Neumann} (Legge di) 144
  \item Ferromagnetismo 241, 249$\div$255
  \item Flusso 10
    \subitem concatenato 174
    \subitem tagliato 174
    \subitem (Tubo di) 14
    \subitem (v.\ Conservazione del flusso \dots)
  \item Forza
    \subitem coercitiva 254
    \subitem di Lorentz 75, 117
    \subitem elettrica su un dipolo 232, 242
    \subitem su un dielettrico 233--234
    \subitem magnetica su una spira 138, 175--176
    \subitem -- su un dipolo magnetico 138, 242
    \subitem su una carica in moto 111, 117
  \item Forza elettromotrice (f.e.m.) 68, 144
    \subitem autoindotta 148
    \subitem indotta 144
  \item Fotone 212
  \item Frange di interferenza 209
  \item \textsc{Franklin} 1
    \subitem (franklin, Fr: Unit\`a) 2
  \item Frequenza 200
    \subitem di \textsc{Larmor} 245
  \item Funzioni armoniche 35
  \item Fusione nucleare 105

  \indexspace
  \item \textbf{G}
  \indexspace
  \item \textsc{Gauss} (gauss, G: Unit\`a) 75
    \subitem (Sistema di) 2--3
    \subitem (Superficie di) 19
    \subitem (Teorema di) 13
  \item Generatore
    \subitem di forza elettromotrice 68
    \subitem di corrente 139
  \item Gradiente 23
  \item \textsc{Green} (Formula di) 36
  \item Gruppo di invarianza delle sorgenti 17
  \item Gruppo di Lorentz 85
    \subitem proprio 85
      \subsubitem -- ortocrono 87, 91

  \indexspace
  \item \textbf{H}
  \indexspace
  \item \textsc{Hafele} e \textsc{Keating} (Esperimento di) 82
  \item \textsc{Hall} (Effetto) 140--141
    \subitem (Resistenza) 140
  \item \textsc{Henry} (henry, H: Unit\`a) 151--152

  \indexspace
  \item \textbf{I}
  \indexspace
  \item Impedenza 171
  \item Induttanza 158
  \item Induzione dielettrica \ $\vec D$ \ 223
  \item Induzione elettrostatica 42
  \item Induzione mutua (v.\ Coefficienti)
  \item Intensit\`a (di un'onda piana) 201
  \item Interferenza 208--209
    \subitem costruttiva 208
    \subitem distruttiva 208
    \subitem (Frange di) 209
  \item Interferometro di \textsc{Mach} -- \textsc{Zehnder} 213
  \item Invarianti
    \subitem del campo elettromagnetico 117
    \subitem (Grandezze) 80, 86
  \item Invarianza di scala 22
  \item Inversione del tempo 142
  \item Isteresi (Ciclo di) 254
    \subitem (Lavoro nel) 259

  \indexspace
  \item \textbf{J}
  \indexspace
  \item \textsc{Joule} (Effetto) 64
    \subitem (joule, J: Unit\`a) 69
    \subitem (Potenza) 64

  \indexspace
  \item \textbf{K}
  \indexspace
  \item \textsc{Keating} (v.\ \textsc{Hafele})
  \item \textsc{Kelvin} (v.\ Effetto)

  \indexspace
  \item \textbf{L}
  \indexspace
  \item Lamina di corrente 115, 134
  \item \textsc{Laplace} (Equazione di) 35
  \item \textsc{Larmor} (Frequenza di) 245
    \subitem (Teorema di) 246
  \item Lavoro di estrazione 41
  \item Lavoro nel ciclo di isteresi 259
  \item Legge
    \subitem di \textsc{Biot} -- \textsc{Savart} 127
    \subitem di densit\`a delle linee di forza 14
    \subitem di \textsc{Faraday} -- \textsc{Neumann} 144
    \subitem -- per i circuiti in movimento 173
    \subitem di \textsc{Lenz} 148
    \subitem di \textsc{Malus} 212--213
    \subitem di \textsc{Ohm} 62, 65--66
    \subitem -- in forma locale 62
    \subitem di trasformazione 78
    \subitem -- del campo elettrico 110, 111, 116
    \subitem -- del campo elettromagnetico 116
    \subitem -- della carica 79--80
    \subitem -- della massa 79--80, 102--103
    \subitem -- delle sorgenti 114
  \item \textsc{Lenz} (Legge di) 148
  \item Linea di universo 92
  \item Linee di forza 5--6
    \subitem fra i conduttori 39
    \subitem (Legge di densit\`a delle) 14
  \item Livelli energetici 212
  \item \textsc{Lorentz} (v.\ anche \textsc{Drude})
    \subitem (Forza di) 75, 117
    \subitem (Gruppo di) 85, 87
    \subitem (Trasformazioni di) 83$\div$88
    \subitem -- senza rotazioni 86$\div$88
  \item Luce polarizzata 213
  \item Lunghezza (o profondit\`a) di penetrazione 155
  \item Lunghezza d'onda 146, 200
  \item Lunghezza propria 90

  \indexspace
  \item \textbf{M}
  \indexspace
  \item \textsc{Mach} -- \textsc{Zehnder} (Interferometro di) 213
  \item Magnete cilindrico 249$\div$251
    \subitem toroidale 251
  \item Magnetizzazione (Curva di) 253
    \subitem (Densit\`a di) 240, 242, 248
    \subitem di saturazione 252, 254
    \subitem in un magnete permanente 252, 253
    \subitem residua 254
  \item Magnetone di \textsc{Bohr} 244
  \item \textsc{Malus} (Legge di) 212--213
  \item Massa di una particella 102
    \subitem di un sistema 104
    \subitem ed energia 102$\div$105
    \subitem nulla (particelle di) 102
  \item Massa (o Terra) 43
  \item \textsc{Millikan} 2
  \item \textsc{Minkowski} (Tensore metrico di) 84, 92
  \item MKSA (Sistema) 2
  \item Mobilit\`a 38, 66
  \item Molecole polari 217, 238--239
    \subitem non polari 217, 238
  \item Momento angolare 244
    \subitem canonico 247
    \subitem cinetico 247
  \item Momento di dipolo, di quadrupolo 26
  \item Momento di dipolo indotto 217--218
    \subitem proprio (o permanente) 217, 219--220
  \item Momento magnetico 136, 243
    \subitem indotto 242, 244$\div$247
    \subitem intrinseco 244
    \subitem orbitale 244
    \subitem proprio (o permanente) 242, 244
  \item Momento su un dipolo
    \subitem elettrico 219
    \subitem magnetico 137
  \item \textsc{Mossotti} -- \textsc{Clausius} (Formula di) 223
  \item Moto di una carica in campo magnetico 100--101
  \item Moto iperbolico 100
  \item Motore 175, 177, 181
    \subitem a induzione 182
  \item Multipoli 27

  \indexspace
  \item \textbf{N}
  \indexspace
  \item \textsc{Neumann} (Problema di) 43

  \indexspace
  \item \textbf{O}
  \indexspace
  \item \textsc{Oersted} (oersted, Oe: unit\`a) 240
  \item \textsc{Ohm} (Legge di) 62, 65--66
    \subitem (ohm, $\Omega$: Unit\`a) 62
  \item Onda
    \subitem (Frequenza di un'o.) 200
    \subitem (Intensit\`a di un'o.) 201
    \subitem (Lunghezza d'o.) 200
    \subitem (Periodo di un'o.) 200
    \subitem (Polarizzazione di un'o.) 198, 201
    \subitem (Propagazione rettilinea di un'o.) 204
    \subitem (Vettore d'o.) 201
  \item Onde (Equazione delle) 196
    \subitem longitudinali 198
    \subitem monocromatiche 200
    \subitem piane 197$\div$199
    \subitem -- progressive 198
    \subitem trasversali 198
  \item Ottupolo 28

  \indexspace
  \item \textbf{P}
  \indexspace
  \item Paramagnetismo 241
  \item Pelle (Effetto) 63, 154--155
  \item Periodo (di un'onda monocromatica) 200
  \item Permalloy 255
  \item Permeabilit\`a magnetica 242
    \subitem del vuoto 126
  \item Permittivit\`a del vuoto 2
  \item Pila a gravit\`a 70
  \item \textsc{Planck} (Costante di) 41, 212, 244
  \item \textsc{Poincar\'e} 195
    \subitem (Campo degli sforzi di) 195
  \item \textsc{Poisson} (Equazione di) 28
  \item Polari (Molecole) 217
  \item Polarizzabilit\`a
    \subitem elettrica 218, 219, 238--239
    \subitem magnetica 247
  \item Polarizzazione (Corrente di) 239
  \item Polarizzazione (Densit\`a di) 217, 238--239
  \item Polarizzazione (di un'onda piana) 198, 201
  \item Polaroid 212
  \item Potenza irraggiata da un'antenna 205
  \item Potenza \textsc{Joule} 64
  \item Potenziale 22
    \subitem (Caduta di) 69
  \item Potenziale magnetico 249
  \item Potenziale vettore 126
    \subitem ritardato 199, 202
  \item \textsc{Poynting} (Teorema di) 189
    \subitem (Vettore di) 189
  \item Problema di \textsc{Dirichlet} 43
    \subitem di \textsc{Neumann} 43
    \subitem misto 43
  \item Profondit\`a (o lunghezza) di penetrazione 155
  \item Propagazione rettilinea di un'onda 204

  \indexspace
  \item \textbf{Q}
  \indexspace
  \item Quadricorrente 114
  \item Quadriforza 105
  \item Quadrivelocit\`a 95
  \item Quadrivettore 83
    \subitem energia--quantit\`a di moto 101
  \item Quadrupolo 26, 28
  \item Quantit\`a di moto 98$\div$100
    \subitem (Conservazione della) 99, 103
    \subitem del campo elettromagnetico 191
    \subitem -- densit\`a 191
    \subitem nascosta 193

  \indexspace
  \item \textbf{R}
  \indexspace
  \item Raggio classico dell'elettrone 41, 194
  \item Raggio di \textsc{Bohr} 41
  \item Rapidit\`a 87
  \item Resistenza 62
    \subitem \textsc{Hall} 140
    \subitem interna 69
  \item Resistivit\`a 62
  \item Riferimento
    \subitem del centro di massa 104
    \subitem tangente 95
  \item Rifrazione delle linee di forza 224
  \item Rotazione (o Rotore) 23, 122--123

  \indexspace
  \item \textbf{S}
  \indexspace
  \item \textsc{Seifert} (Superficie di) 124
  \item Sfera dielettrica in campo esterno 226--227
  \item Sfera magnetizzata 255
  \item Simultaneit\`a 90
  \item Sistema di unit\`a 2
  \item Solenoide 129$\div$131, 134$\div$135, 147$\div$150
  \item Sostanze
    \subitem diamagnetiche 241
    \subitem ferromagnetiche 241, 249$\div$255
    \subitem paramagnetiche 241
    \subitem polari 217
  \item Spire superconduttrici (conservazione del flusso nelle) 188
  \item Stabilizzatore 17
  \item \textsc{Stokes} (Teorema di) 123
  \item Superconduttivit\`a 67
  \item Superficie di \textsc{Seifert} 124
  \item Supermalloy 255
  \item Suscettivit\`a
    \subitem elettrica 221
    \subitem magnetica 242, 247

  \indexspace
  \item \textbf{T}
  \indexspace
  \item Temperatura di \textsc{Curie} 241
  \item Tempo proprio 89, 95
  \item Tensione 149
  \item Teorema
    \subitem della divergenza (o di \textsc{Green}) 12
    \subitem della conservazione dell'energia 103
    \subitem della media 35
    \subitem delle forze vive 101
    \subitem di \textsc{Amp\`ere} 125
    \subitem di \textsc{Gauss} 13
    \subitem di \textsc{Larmor} 246
    \subitem di \textsc{Poynting} 189
    \subitem di reciprocit\`a 151
    \subitem di \textsc{Stokes} 123
  \item Terra (o Massa) 43
  \item \textsc{Tesla} (tesla, T: Unit\`a) 75
  \item Terzo principio della dinamica relativistica 97--98
  \item Trasformatore 160--161

  \indexspace
  \item \textbf{V}
  \indexspace
  \item Velocit\`a
    \subitem (Composizione delle) 88, 90--91, 94$\div$96
    \subitem delle cariche in un conduttore 60, 66
    \subitem quadratica media 60
  \item Vettore
    \subitem di \textsc{Poynting} 189
    \subitem d'onda 201
    \subitem isotropo (o tipo luce) 91, 92
    \subitem tipo spazio, tipo tempo 91
  \item Vita media 82

  \indexspace
  \item \textbf{W}
  \indexspace
  \item \textsc{Weber} (weber, Wb: Unit\`a) 75
  \item \textsc{Weiss} (Domini di) 252

  \indexspace
  \item \textbf{Z}
  \indexspace
  \item \textsc{Zehnder} (v.\ \textsc{Mach} -- \textsc{Zehnder})

\end{theindex}
